\section{Theory Task T3: Optimization Interface and Conditional Convergence}
\label{sec:theory-t3-optimization}

T1 freezes the Event-FedAvg transition and T2 characterizes its stateful
encoder.  This section couples the emitted pulses back to a smooth global
objective.  The result is deliberately conditional: it identifies the exact
aggregate alignment quantity that must be controlled for convergence, while
using T2 to control the event-energy and communication terms.  It does not
reintroduce either of the descent certificates investigated during mechanism
development.

\begin{takeawaybox}[title={T3 result boundary}]
For an $L$-smooth, lower-bounded objective, Event-FedAvg obeys an exact
descent-versus-harm decomposition.  Under an explicit conditional aggregate
alignment assumption, it has a finite-horizon stationarity bound, a
constant-quantum neighborhood bound, and a diminishing-quantum asymptotic
corollary.  T2 bounds the curvature penalty by the realized coordinate-event
stream.  What remains unproved is the alignment assumption itself under
general heterogeneous, model-coupled FedAvg evidence.
\end{takeawaybox}

\subsection{Round-level quantities}

Let
\begin{equation}
    F(w)=\sum_{i=1}^{M}p_iF_i(w),
    \qquad p_i>0,
    \qquad \sum_{i=1}^{M}p_i=1,
    \label{eq:t3-global-objective}
\end{equation}
and write $g^r=\nabla F(w^r)$.  Under the frozen full-participation semantics,
client $i$ emits $c_i^r\in\{-1,0,+1\}^d$, the server forms
\begin{equation}
    C^r=\sum_{i=1}^{M}c_i^r,
    \label{eq:t3-aggregate-pulse}
\end{equation}
and applies
\begin{equation}
    w^{r+1}=w^r+q_rC^r,
    \qquad q_r>0.
    \label{eq:t3-server-update}
\end{equation}
Define the number of client-coordinate events in round $r$ by
\begin{equation}
    N_r=\sum_{i=1}^{M}\sum_{j=1}^{d}(c_{ij}^r)^2
       =\sum_{i=1}^{M}\|c_i^r\|_0.
    \label{eq:t3-round-event-count}
\end{equation}
Thus $0\leq N_r\leq Md$.  This counts client-coordinate events before
same-coordinate pulses are added at the server; it is the quantity that also
drives the sparse uplink accounting in T2.

For each realized round, separate the gradient mass carried by pulses that
oppose the global gradient from the mass carried by pulses that follow it:
\begin{align}
    D_r
    &=
    \sum_{i=1}^{M}\sum_{j=1}^{d}
    |g_j^r|\,\mathbf{1}_{\{c_{ij}^rg_j^r<0\}},
    \label{eq:t3-descent-mass}\\
    H_r
    &=
    \sum_{i=1}^{M}\sum_{j=1}^{d}
    |g_j^r|\,\mathbf{1}_{\{c_{ij}^rg_j^r>0\}}.
    \label{eq:t3-harmful-mass}
\end{align}
The letters stand for descent-aligned and harmful gradient mass.  A silent
coordinate, or one with $g_j^r=0$, contributes to neither quantity.

\begin{proposition}[Exact aggregate alignment identity]
\label{prop:t3-alignment-identity}
For every realized round,
\begin{equation}
    -\langle g^r,C^r\rangle=D_r-H_r.
    \label{eq:t3-alignment-identity}
\end{equation}
This identity already includes cancellation among clients that emit opposing
pulses on the same coordinate.
\end{proposition}

\begin{proof}
Expanding the inner product gives
\begin{equation}
    \langle g^r,C^r\rangle
    =\sum_{i=1}^{M}\sum_{j=1}^{d}g_j^rc_{ij}^r.
\end{equation}
Each nonzero summand is $-|g_j^r|$ when the pulse opposes the gradient and
$+|g_j^r|$ when it follows the gradient.  Summing the two cases gives
\eqref{eq:t3-alignment-identity}.
\end{proof}

\begin{lemma}[Aggregate event energy]
\label{lem:t3-event-energy}
For every round,
\begin{equation}
    \|C^r\|_2^2\leq MN_r\leq M^2d.
    \label{eq:t3-event-energy}
\end{equation}
Consequently, for any horizon $R$ and $q_{\max,R}=\max_{1\leq r\leq R}q_r$,
\begin{align}
    \sum_{r=1}^{R}q_r^2\|C^r\|_2^2
    &\leq
    M\sum_{r=1}^{R}q_r^2N_r
    \label{eq:t3-weighted-event-energy}\\
    &\leq
    Mq_{\max,R}^2N_R^{\mathrm{coord}}.
    \label{eq:t3-event-energy-cumulative}
\end{align}
With a common threshold $\vartheta$, T2 further gives
\begin{equation}
    \sum_{r=1}^{R}q_r^2\|C^r\|_2^2
    \leq
    \frac{Mq_{\max,R}^2}{\vartheta}
    \left(
      \sum_{i=1}^{M}\|z_i^1\|_1
      +\sum_{r=1}^{R}\sum_{i=1}^{M}\|a_i^r\|_1
    \right).
    \label{eq:t3-energy-evidence-budget}
\end{equation}
\end{lemma}

\begin{proof}
For each coordinate, Cauchy--Schwarz gives
\begin{equation}
    \left(\sum_{i=1}^{M}c_{ij}^r\right)^2
    \leq M\sum_{i=1}^{M}(c_{ij}^r)^2.
\end{equation}
Summing over $j$ proves the first inequality.  The remaining statements follow
from $N_r\leq Md$, from
$\sum_rN_r=N_R^{\mathrm{coord}}$, and from
Proposition~\ref{prop:t2-vector-event-budget}.
\end{proof}

The factor $M$ is a worst-case coherence factor.  Opposing simultaneous pulses
can make $\|C^r\|_2^2$ much smaller than $MN_r$, but they do not reduce the
uplink event count $N_r$.

\subsection{A deterministic descent inequality}

\begin{proposition}[One-step Event-FedAvg descent decomposition]
\label{prop:t3-one-step-descent}
Suppose $F$ has an $L$-Lipschitz gradient.  Then every realized Event-FedAvg
round satisfies
\begin{equation}
    F(w^{r+1})
    \leq
    F(w^r)
    -q_r(D_r-H_r)
    +\frac{LM}{2}q_r^2N_r.
    \label{eq:t3-one-step-descent}
\end{equation}
\end{proposition}

\begin{proof}
The smoothness inequality and \eqref{eq:t3-server-update} give
\begin{equation}
    F(w^{r+1})
    \leq
    F(w^r)
    +q_r\langle g^r,C^r\rangle
    +\frac{L}{2}q_r^2\|C^r\|_2^2.
\end{equation}
Apply Proposition~\ref{prop:t3-alignment-identity} and
Lemma~\ref{lem:t3-event-energy}.
\end{proof}

Equation~\eqref{eq:t3-one-step-descent} is the optimization interface supplied
by T3.  The net first-order benefit is $D_r-H_r$; the second-order cost is
controlled by the model quantum and the realized event stream.  Silence is not
automatically beneficial: when $N_r=D_r=H_r=0$, the model simply does not move.

\subsection{Conditional finite-horizon stationarity}

Let $\mathcal{F}_r$ contain the initial state and all randomness revealed before
the local computation of round $r$.  In particular, $w^r$ and $g^r$ are
$\mathcal{F}_r$-measurable.  The following condition states exactly what is
needed from the closed encoder--optimization loop.

\paragraph{Aggregate alignment with defect.}
There exist a constant $\kappa>0$ and deterministic numbers $\beta_r\geq0$
such that
\begin{equation}
    \mathbb{E}[D_r-H_r\mid\mathcal{F}_r]
    \geq
    \kappa\|\nabla F(w^r)\|_2^2-\beta_r.
    \label{eq:t3-alignment-assumption}
\end{equation}
The defect $\beta_r$ is not an algorithmic safeguard.  It records the loss of
global descent alignment caused by silence, leakage, stale accumulated
evidence, local-update drift, stochastic sign error, and client heterogeneity.

\begin{proposition}[Conditional finite-horizon stationarity bound]
\label{prop:t3-finite-horizon}
Assume that $F$ is $L$-smooth and lower bounded by $F_{\inf}$, and that
\eqref{eq:t3-alignment-assumption} holds.  Then, for every $R\geq1$,
\begin{equation}
    \kappa\sum_{r=1}^{R}q_r
    \mathbb{E}\|\nabla F(w^r)\|_2^2
    \leq
    F(w^1)-F_{\inf}
    +\sum_{r=1}^{R}q_r\beta_r
    +\frac{LM}{2}\sum_{r=1}^{R}q_r^2\mathbb{E}N_r.
    \label{eq:t3-finite-horizon-bound}
\end{equation}
Let $Q_R=\sum_{r=1}^{R}q_r$ and draw $J_R\in\{1,\ldots,R\}$ independently
with $\mathbb{P}(J_R=r)=q_r/Q_R$.  Then
\begin{equation}
    \mathbb{E}\|\nabla F(w^{J_R})\|_2^2
    \leq
    \frac{
      F(w^1)-F_{\inf}
      +\sum_{r=1}^{R}q_r\beta_r
      +\frac{LM}{2}\sum_{r=1}^{R}q_r^2\mathbb{E}N_r
    }{\kappa Q_R}.
    \label{eq:t3-random-iterate-bound}
\end{equation}
The same right-hand side bounds
$\min_{1\leq r\leq R}\mathbb{E}\|\nabla F(w^r)\|_2^2$.
\end{proposition}

\begin{proof}
Take conditional expectation in \eqref{eq:t3-one-step-descent} and apply
\eqref{eq:t3-alignment-assumption}.  Taking total expectation and summing over
$r$ telescopes the objective values.  Lower bounding the final objective by
$F_{\inf}$ proves \eqref{eq:t3-finite-horizon-bound}.  Division by $\kappa Q_R$
gives the random-iterate identity and the minimum bound.
\end{proof}

\begin{corollary}[Constant model quantum]
\label{cor:t3-constant-quantum}
Let $q_r=q$, and suppose
\begin{equation}
    \frac{1}{R}\sum_{r=1}^{R}\beta_r\leq\bar\beta_R,
    \qquad
    \frac{1}{R}\sum_{r=1}^{R}\mathbb{E}N_r\leq\bar N_R.
\end{equation}
Then
\begin{equation}
    \min_{1\leq r\leq R}
    \mathbb{E}\|\nabla F(w^r)\|_2^2
    \leq
    \frac{F(w^1)-F_{\inf}}{\kappa qR}
    +\frac{\bar\beta_R}{\kappa}
    +\frac{LM\bar N_R}{2\kappa}q.
    \label{eq:t3-constant-quantum-bound}
\end{equation}
In the worst case, $\bar N_R\leq Md$.
\end{corollary}

This exposes two distinct neighborhood terms.  Decreasing $q$ reduces the
smoothness penalty but does not by itself reduce the alignment defect.  In
particular, a leaky encoder that remains silent near a nonstationary point must
pay for that silence through $\beta_r$.

\begin{corollary}[Diminishing model quantum]
\label{cor:t3-diminishing-quantum}
Under the assumptions of Proposition~\ref{prop:t3-finite-horizon}, suppose
\begin{equation}
    \sum_{r=1}^{\infty}q_r=\infty,
    \qquad
    \sum_{r=1}^{\infty}q_r^2<\infty,
    \qquad
    \sum_{r=1}^{\infty}q_r\beta_r<\infty.
    \label{eq:t3-robbins-monro-window}
\end{equation}
Then
\begin{equation}
    \lim_{R\to\infty}
    \frac{1}{Q_R}
    \sum_{r=1}^{R}q_r
    \mathbb{E}\|\nabla F(w^r)\|_2^2
    =0,
    \label{eq:t3-weighted-stationarity}
\end{equation}
and therefore
\begin{equation}
    \liminf_{r\to\infty}
    \mathbb{E}\|\nabla F(w^r)\|_2^2=0.
    \label{eq:t3-liminf-stationarity}
\end{equation}
\end{corollary}

\begin{proof}
Use $N_r\leq Md$ in \eqref{eq:t3-finite-horizon-bound}.  The numerator of the
weighted-average bound remains finite while $Q_R$ diverges.  A nonnegative
sequence whose weighted average tends to zero with positive divergent total
weight must have limit inferior zero.
\end{proof}

This corollary establishes weighted stationarity and a limit-inferior result,
not convergence of every iterate.  Almost-sure convergence or convergence of
the full gradient sequence would require additional martingale, stability, or
geometric assumptions.

\subsection{Optional stronger geometry: a PL neighborhood}

\begin{corollary}[Polyak--\L{}ojasiewicz objective]
\label{cor:t3-pl}
Suppose, in addition, that
\begin{equation}
    \|\nabla F(w)\|_2^2
    \geq 2\mu_{\mathrm{PL}}(F(w)-F_\star)
    \label{eq:t3-pl-condition}
\end{equation}
for all visited iterates, with $\mu_{\mathrm{PL}}>0$.  If $q_r=q$,
$0<2\kappa\mu_{\mathrm{PL}}q\leq1$, $\beta_r\leq\bar\beta$, and
$\mathbb{E}N_r\leq\bar N$, then
\begin{align}
    \mathbb{E}[F(w^{R+1})-F_\star]
    \leq{}&
    (1-2\kappa\mu_{\mathrm{PL}}q)^R
    (F(w^1)-F_\star)
    \nonumber\\
    &+\frac{\bar\beta}{2\kappa\mu_{\mathrm{PL}}}
    +\frac{LMq\bar N}{4\kappa\mu_{\mathrm{PL}}}.
    \label{eq:t3-pl-neighborhood}
\end{align}
\end{corollary}

\begin{proof}
Combine the conditional form of \eqref{eq:t3-one-step-descent} with
\eqref{eq:t3-alignment-assumption} and \eqref{eq:t3-pl-condition}.  The
resulting affine recursion has contraction factor
$1-2\kappa\mu_{\mathrm{PL}}q$; summing its geometric series gives the bound.
\end{proof}

The PL condition is not asserted for the neural networks in the empirical
campaign.  This corollary only records what the same Event-FedAvg interface
implies when stronger objective geometry is independently available.

\subsection{How T2 sign reliability enters the alignment defect}

T2 controls the sign of a first-passage event relative to a locally persistent
mean evidence direction.  T3 requires alignment relative to the current global
gradient.  The following elementary split shows the gap between the two.

\begin{lemma}[Local sign error versus global mismatch]
\label{lem:t3-local-global-split}
Fix a client-coordinate pulse $c_{ij}^r$, a reference local mean evidence
$\mu_{ij}^r$, and a nonzero global-gradient coordinate $g_j^r$.  Then
\begin{equation}
    \mathbf{1}_{\{c_{ij}^rg_j^r>0\}}
    \leq
    \mathbf{1}_{\{c_{ij}^r\neq0,\ c_{ij}^r\mu_{ij}^r\leq0\}}
    +
    \mathbf{1}_{\{c_{ij}^r\neq0,\ \mu_{ij}^rg_j^r>0\}}.
    \label{eq:t3-local-global-split}
\end{equation}
\end{lemma}

\begin{proof}
For a harmful nonzero pulse, either its sign fails to match the reference local
mean, including the zero-mean case, or its sign matches that mean.  In the
latter case the local mean has the same sign as the global gradient, so the
second indicator equals one.
\end{proof}

The first term is an encoder sign-error term and can be bounded by T2 on blocks
that satisfy its stationary sub-Gaussian assumptions.  The second is a
local-to-global mismatch term caused by client heterogeneity, local-update
drift, or stale accumulation; it is not controlled by the encoder theorem.
Thus T2 sign reliability is useful but insufficient on its own to establish
\eqref{eq:t3-alignment-assumption}.

\subsection{Schedule audit for the frozen experiments}

The final matched-baseline campaign uses
\begin{equation}
    q_r=q_0\left(1+\frac{r}{r_0}\right)^{-\alpha}.
    \label{eq:t3-power-law-schedule}
\end{equation}
For an infinite run, the two step-size series in
\eqref{eq:t3-robbins-monro-window} satisfy
\begin{align}
    \sum_rq_r=\infty
    &\quad\Longleftrightarrow\quad \alpha\leq1,
    \\
    \sum_rq_r^2<\infty
    &\quad\Longleftrightarrow\quad \alpha>\frac12.
    \label{eq:t3-power-law-window}
\end{align}
Hence the classical asymptotic window is
\begin{equation}
    \frac12<\alpha\leq1.
    \label{eq:t3-alpha-window}
\end{equation}

\begin{center}
\begin{tabular}{lcccccc}
\toprule
Model & $q_0$ & $r_0$ & $\alpha$ & rounds & $\sum q_r$ & $\sum q_r^2$\\
\midrule
MLP & $0.005$ & $100$ & $0.1$ & $150$ & diverges & diverges\\
CNN & $0.010$ & $100$ & $0.3$ & $80$ & diverges & diverges\\
T3 asymptotic family & free & $>0$ & $(1/2,1]$ & infinite & diverges & converges\\
\bottomrule
\end{tabular}
\end{center}

The first two rows describe finite experimental horizons, for which both sums
are finite numerically.  The labels in the final two columns describe what
would happen if the same formulas were extended indefinitely.  Therefore the
reported MLP and CNN schedules are finite-horizon empirical schedules, not
instances of Corollary~\ref{cor:t3-diminishing-quantum}.  Establishing an
asymptotic result requires either a schedule in \eqref{eq:t3-alpha-window} or a
different argument that exploits stronger decay in the realized event stream.

\subsection{A verifiable offline alignment audit}

The theory assumption can be measured without changing the training rule.  On
selected rounds, store the client pulse vectors and evaluate a sufficiently
accurate reference global gradient.  Then compute
\begin{align}
    A_r^{\mathrm{net}}
    &=-\langle\nabla F(w^r),C^r\rangle=D_r-H_r,
    \label{eq:t3-audit-net-alignment}\\
    \chi_r
    &=\frac{H_r}{D_r+H_r}
      \quad\text{when }D_r+H_r>0,
    \label{eq:t3-audit-harm-fraction}\\
    \beta_r(\kappa)
    &=\left[
       \kappa\|\nabla F(w^r)\|_2^2-A_r^{\mathrm{net}}
      \right]_+.
    \label{eq:t3-audit-defect}
\end{align}
Together with $N_r$, these values separate net alignment, harmful-event share,
alignment defect, and event energy.  They are diagnostics, not acceptance
tests: no event is blocked and no extra server update is applied.

The existing campaign records accuracy, objective, and communication totals,
but it was not designed to retain the full per-round pulse/gradient information
needed for this audit.  A targeted instrumented rerun is therefore the direct
empirical test of T3's only method-specific convergence assumption.

\subsection{T3 conclusions and remaining theorem gap}

T3 establishes the following scoped result:
\begin{enumerate}
    \item every aggregate pulse has an exact descent-aligned minus harmful
    gradient-mass decomposition;
    \item smoothness converts that decomposition into a deterministic
    one-step objective bound;
    \item the quadratic remainder is bounded by the actual coordinate-event
    count and, through T2, by accumulated evidence variation;
    \item explicit aggregate alignment with summable defect yields a
    finite-horizon stationarity bound and a diminishing-quantum asymptotic
    corollary;
    \item constant $q$ produces separate alignment-defect and event-energy
    neighborhoods;
    \item local first-passage sign reliability does not eliminate the distinct
    client-to-global mismatch caused by heterogeneity and local drift;
    \item the frozen experimental schedules support finite-horizon claims but
    do not satisfy the stated asymptotic square-summability condition.
\end{enumerate}

The unresolved theoretical task is now precise: bound or empirically calibrate
$\beta_r$ for model-coupled multi-step FedAvg evidence.  A next proof should
decompose that quantity into stochastic encoder sign error, local-SGD drift,
heterogeneity, stale-state error, and silence/deadzone terms.  Until such a
bound is available, Proposition~\ref{prop:t3-finite-horizon} is a conditional
convergence interface rather than an unconditional convergence theorem for the
complete frozen method.
